%% Dans ce fichier, les �ventuels accents sont cod�s
%% en ISO-8859-15 (aka latin9)
%% La ligne suivante donne trois exemples 
%% e accent aigu : �, e accent grave : �, e-dans-l'o : �
%%%%%ENTREES%%%%
\begin{tabarg}
\multicolumn{3}{l}{\titretab {Entr\'ees}}\\
\hline
$JDA$ & Coeff., \texttt{int} & Num�ro du jour dans l'ann�e (de 1 � 366)\\
$Time$ & Vecteur, \texttt{double} & Heure UTC (sec. depuis 0h00 du jour pass� en premier argument)\\
$Lon$ & Vecteur, \texttt{float} & Longitude (degree)\\
$Lat$ & Vecteur, \texttt{float} & Latitude (degree)\\
$\psi$ & Vecteur, \texttt{float} & Azimut (ou cap) de la plate-forme de mesure (degree)\\
$\phi_p$ & Vecteur, \texttt{float} & Angle de g�te, ou roulis (degree)\\
$\theta_p$ & Vecteur, \texttt{float} & Assiette longitudinale, ou tangage (degree)\\
$\Delta\phi$ & Coeff., \texttt{double} & D�calage entre l'angle de g�te de l'instrument et celui de la plate-forme (degree)\\
$\Delta\theta$ & Coeff., \texttt{double} & D�calage entre l'angle de tangage de l'instrument et celui de la plate-forme (degree)\\
SW$_d$ & Vecteur, \texttt{float} & Flux global incident (W/m$^2$)\\
\end{tabarg}

\medskip
%%%%%SORTIES%%%%
\begin{tabarg}
\multicolumn{3}{l}{\titretab {Sorties}}\\
\hline
SW$_c$ & Vecteur, \texttt{float} & Flux global incident corrig� (W/m$^2$)\\
\end{tabarg}
